\section{Biological and Psychological Requirements}

\subsection{Human Vision}

ZONES OF VISION

The peripheral vision is especially important to recognize hazards quickly, and should not be obstructed \parencite{}.

How much of the field of view is obstructed, depends on the design and distance of the suit's face shield (see figure below).

GRAPHIC FOV, wide vs narrow face shield, long vs short distance
- wider/less distance to face = bigger view angle

focus on (pre-)peripheral as limit/"frame" for projected UI

what about turning the head?

---

EYE STRAIN with light-emitting displays
does that get worse if the display is closer to the eye?
is eye strain linked to earlier exhaustion/other relevant negative effects?

ui design possibilities to reduce eye strain

trade-offs: low contrast vs visibility?

---

how does acute stress/panic affect vision? what about physical emergencies, like dizziness? fov narrower? vision loss? blur?

\subsection{Cognitive Load}

cognitive load reduced when active?

cognitive offloading

\subsection{Perceptual Psychology}

reaction to color + movement

blinking bad? frequency?

---

speed of comprehending/parsing information

---

- mitigating autopilot mode for routine tasks
